% ============================================================
% KAPITEL 2 — GRUNDLAGEN UND STAND DER TECHNIK
% ============================================================

\chapter{Grundlagen und Stand der Technik}
\label{ch:grundlagen}

Dieses Kapitel führt zunächst die eingesetzten Modelle, das Agenten-Paradigma und die Formen der Szenenrepräsentation ein und ordnet die Arbeit anschließend in den Stand der Technik ein.


% ------------------------------------------------------------
\section{Große Sprach- und Sprach-Bild-Modelle}
\label{sec:grundlagen-modelle}

Die Arbeit setzt zwei Arten vortrainierter Modelle als Bausteine ein, die unverändert genutzt und nicht weiterentwickelt werden.

Ein großes Sprachmodell, im Folgenden als LLM (Large Language Model) bezeichnet, ist ein auf umfangreichen Textmengen trainiertes neuronales Modell, das zu einer Eingabe fortlaufend Text erzeugt.
Durch zusätzliches Training auf Instruktionen lässt es sich anweisen, eine Aufgabe zu bearbeiten, und gibt seine Antwort wiederum als Text aus.
In dieser Arbeit übernimmt ein LLM die Rolle des Agenten, der einen Befehl interpretiert und die passende Operation auswählt.
% CITATION: Grundlagenreferenz zu LLMs bzw. Instruction-Tuning ergaenzen (Survey/Lehrbuch).

Ein Sprach-Bild-Modell, im Folgenden als VLM (Vision-Language Model) bezeichnet, erweitert dieses Prinzip um eine visuelle Eingabe.
Es verarbeitet ein Bild zusammen mit einem Text und erzeugt daraus eine textuelle Antwort, etwa eine Beschreibung des Bildinhalts oder die Antwort auf eine Frage zum Bild.
In dieser Arbeit übernimmt ein VLM die visuelle Wahrnehmung der Szene.
% CITATION: Grundlagenreferenz zu VLMs ergaenzen (Survey).

Die beiden Modelle sind im hier untersuchten System nicht zu einem multimodalen Modell verschmolzen, sondern hintereinandergeschaltet.
Das Sprach-Bild-Modell erhält das gerenderte Bild und erzeugt daraus eine natürlichsprachliche Beschreibung.
Erst diese Beschreibung gelangt in den Kontext des Sprachmodells, das den Befehl deutet und die Operation wählt.
Das Modell, das die eigentliche Entscheidung trifft, sieht also zu keinem Zeitpunkt ein Bild.

Diese Eigenschaft des Aufbaus ist für die Deutung aller späteren Ergebnisse wichtig.
Was in dieser Arbeit \enquote{visuelle Repräsentation} heißt, ist streng genommen keine Bildeingabe, sondern die sprachliche Beschreibung eines Bildes.
Ist später von der Gewichtung der strukturellen gegenüber der visuellen Schicht die Rede, konkurrieren daher nicht Bild und Text im Inneren eines multimodalen Modells, sondern zwei Textquellen im Kontext eines reinen Sprachmodells: eine exakte, aber unvollständige Datenstruktur und eine reichhaltige, aber unscharfe Beschreibung.
Die konkrete Umsetzung zeigt Abschnitt~\ref{sec:impl-variants}, die Folgen für die Befunde diskutiert Abschnitt~\ref{sec:eval-validity}.

Beide Modelle werden lokal im Sinne der Definition aus Abschnitt~\ref{sec:ein-motivation} betrieben und haben eine Größe von jeweils etwa sieben Milliarden Parametern.
Die Gründe für den lokalen Betrieb nennt Abschnitt~\ref{sec:ein-relevanz}.


% ------------------------------------------------------------
\section{Werkzeugnutzende Agenten}
\label{sec:grundlagen-agenten}

Die Interaktion mit der Szene folgt dem Muster eines werkzeugnutzenden Agenten.
Darunter versteht man ein Sprachmodell, das nicht nur Text erzeugt, sondern in einer Schleife aus Wahrnehmung und Handlung Werkzeuge aufruft, um eine Aufgabe zu lösen.

Ein Werkzeug ist eine benannte Funktion mit einer kurzen Beschreibung ihrer Wirkung und ihrer Parameter.
Anhand dieser Beschreibung entscheidet das Sprachmodell, welches Werkzeug es mit welchen Argumenten aufruft.
Das Ergebnis des Aufrufs fließt in den Kontext zurück, und das Modell wählt den nächsten Schritt, bis es die Aufgabe als abgeschlossen meldet.
So verbindet der Agent die sprachliche Interpretation eines Befehls mit konkreten Aktionen in der Umgebung.

Umgesetzt wird der Agent mit dem Framework smolagents, das einen werkzeugaufrufenden Agenten bereitstellt~\cite{smolagents}.
Die Werkzeuge sind über das Model Context Protocol (MCP) angebunden, einen offenen Standard, über den Sprachmodelle auf externe Werkzeuge und Datenquellen über eine einheitliche Schnittstelle zugreifen~\cite{mcp}.
Das Protokoll trennt die Bereitstellung der Werkzeuge von ihrer Nutzung, sodass dieselben Werkzeuge unabhängig vom konkreten Modell aufgerufen werden können.
Wie das System dieses Muster umsetzt, beschreibt Kapitel~\ref{ch:implementierung}.


% ------------------------------------------------------------
\section[Szenenrepräsentation]{Szenenrepräsentation für Sprachmodelle}
\label{sec:grundlagen-repraesentation}

Damit ein Sprachmodell eine Szene verändern kann, muss es sie zunächst erfassen, und da es auf Text arbeitet, muss die dreidimensionale Szene in eine dafür zugängliche Form gebracht werden.
Dafür gibt es zwei grundlegende Ansätze.

Der strukturelle Ansatz beschreibt die Szene als Daten.
Die Objekte werden mit Eigenschaften wie Typ, Position, Größe und Farbe in textueller Form aufgeführt, etwa als JSON oder als Szenengraph.
Diese Darstellung ist exakt und für ein Sprachmodell unmittelbar lesbar, gibt aber nur wieder, was in den Daten kodiert ist.

Der visuelle Ansatz stellt die Szene als gerendertes Bild bereit, das ein Sprach-Bild-Modell interpretiert.
Er erfasst das Aussehen der Szene, wie es sich einem Betrachter zeigt, liefert aber keine exakten Koordinaten.
Hinzu kommt, dass ein Modell ein im Bild erkanntes Objekt benennen können muss, um es danach gezielt anzusprechen.
Verbreitet ist dafür das Set-of-Mark-Prompting, bei dem sichtbare Objekte im Bild mit nummerierten Marken versehen werden~\cite{yang2023}.
Das Modell kann sich dann auf eine Marke beziehen, deren Nummer sich eindeutig einem Objekt zuordnen lässt.

Wie schwierig eine Aufgabe ist, hängt unabhängig von der Repräsentationsform auch davon ab, wie der Befehl auf ein Objekt verweist.
Bestehende Arbeiten unterscheiden Referenzen etwa danach, ob sie vom Blickwinkel des Betrachters abhängen oder ob das Zielobjekt in der Szene eindeutig ist~\cite{achlioptas2020, chen2020}.
Von solchen Einteilungen geht die feinere Aufschlüsselung nach Referenztypen aus, die Abschnitt~\ref{sec:eval-taxonomy} einführt.


% ------------------------------------------------------------
\section{Stand der Technik}
\label{sec:grundlagen-sota}

Die einzelnen Bausteine der Arbeit sind gut erforscht.
Dieser Abschnitt ordnet die Arbeit in vier angrenzende Forschungsfelder ein und zeigt jeweils, wo sich deren Fragestellung von der hier untersuchten unterscheidet.

\subsection{Strukturelle Szenenmanipulation in immersiven Umgebungen}

Ein erstes Feld befasst sich mit der sprachgesteuerten Manipulation von Szenen in virtuellen und gemischten Realitäten.
LLMR erzeugt und verändert interaktive Mixed-Reality-Szenen zur Laufzeit, indem es mehrere Sprachmodelle orchestriert und die Szene als strukturierten Graphen führt~\cite{delatorre2024}.
VR Mover manipuliert Objekte in der virtuellen Realität über natürliche Sprache und kodiert die Szene dazu als strukturierte Objektdaten~\cite{wang2025}.
Beide Arbeiten zeigen, dass sich Szenen in immersiven Umgebungen sprachlich verändern lassen, nutzen dafür aber ausschließlich die strukturelle Repräsentation und große, cloudbasierte Modelle.
Verschiedene Repräsentationsformen vergleichen sie nicht.

\subsection{Hybrides Grounding}

Ein zweites Feld verbindet strukturelle und visuelle Information, um ein Objekt in einer Szene zu bestimmen.
SeeGround stellt eine Szene als Kombination aus einer gerenderten Ansicht und einer räumlich angereicherten textuellen Beschreibung dar und verknüpft beide über visuelle Marken~\cite{li2025}.
Konzeptionell steht dieser Ansatz der vorliegenden Arbeit am nächsten, verfolgt aber ein anderes Ziel: Er dient dem Auffinden eines Objekts, nicht seiner Manipulation, setzt ein großes Modell ein und stellt die Repräsentationsformen nicht als gleichwertige Bedingungen einander gegenüber.

Ähnlich geht SceneGraphGrounder vor, das das Auffinden eines Objekts als strukturiertes Graph-Matching formuliert~\cite{sun2026}.
Ein visuelles Modell leitet dabei über eine Marker-Strategie die Beziehungen zwischen Objekten aus zweidimensionalen Ansichten ab, die anschließend in einen persistenten Szenengraphen überführt werden.
Die Autoren berichten, dass eine explizite strukturelle Repräsentation die Grounding-Genauigkeit verbessert, was zeigt, dass die Verbindung beider Informationsarten aktiv verfolgt wird.
Wie SeeGround bleibt die Arbeit aber beim Auffinden von Objekten und setzt strukturelle und visuelle Information von vornherein gemeinsam ein, statt sie als getrennte Repräsentationsformen gegeneinander zu prüfen.

\subsection{Modalitäts-Ablationen und Modalitäts-Dominanz}

Ein drittes Feld untersucht, wie sich verschiedene Eingabemodalitäten auf die Leistung auswirken.
Die Arbeit zum Embodied Agent Interface vergleicht unter anderem eine rein strukturelle, eine rein visuelle und eine kombinierte Eingabe und zeigt, dass strukturierter Input reine Bildeingabe deutlich übertrifft~\cite{li2025a}.
Der Vergleich betrifft allerdings die Handlungsplanung, nicht die Szenenmanipulation, und schlüsselt die Aufgaben nicht nach Referenztypen auf.
In eine ähnliche Richtung weist SIRI-Bench, ein Benchmark für räumliches Schließen~\cite{song2026}.
In einer kontrollierten Ablation werden dort baugleiche Modelle einmal mit textueller und einmal mit visueller Eingabe derselben Aufgabe verglichen.
Die textbasierten Varianten übertreffen ihre videobasierten Gegenstücke deutlich; die Modelle stützen sich übermäßig auf textuelle Hinweise und gewinnen räumliche Information nur schwer aus dem visuellen Eingang.
Das stützt aus einer anderen Domäne die auch in dieser Arbeit beobachtete Tendenz, dass strukturierte Information bei räumlichen Aufgaben gegenüber rein visueller im Vorteil ist, betrifft aber abstrakte Reasoning-Aufgaben statt der Manipulation einer 3D-Szene und unterscheidet keine Referenztypen.

Eng verwandt sind zwei Arbeiten zur Modalitäts-Dominanz.
Sie zeigen, dass Sprach-Bild-Modelle textuelle gegenüber visueller Information übergewichten und bei widersprüchlichen Eingaben eher dem Text vertrauen~\cite{deng2025, wu2025}.
% [INFO BENOETIGT] Bitte in wu2025 die Liste der evaluierten Modelle
% gegenpruefen und den folgenden Satz entsprechend praezisieren oder
% streichen.
Eine der beiden Arbeiten untersucht dabei auch ein Modell derselben Familie, aus der das hier eingesetzte visuelle Modell stammt~\cite{wu2025}.

Auf die vorliegende Untersuchung lassen sich diese Befunde allerdings nur als Analogie übertragen.
Beide Arbeiten betrachten die Gewichtung von Bild gegen Text \emph{innerhalb} eines multimodalen Modells, dem beide Modalitäten zugleich vorliegen.
In der hier verwendeten zweistufigen Pipeline sieht das entscheidende Sprachmodell dagegen gar kein Bild, sondern zwei Textquellen (Abschnitt~\ref{sec:grundlagen-modelle}).
Der Konflikt besteht also zwischen einer präzisen, maschinell erzeugten Datenstruktur und einer sprachlich vagen Schilderung derselben Szene.
Eine direkte Erwartung für diesen Aufbau begründen die Arbeiten daher nicht; sie legen aber als strukturell ähnliches Muster nahe, dass eine naive Kombination beider Quellen die präzisere von beiden überbewerten könnte.
Wo sich spätere Kapitel auf diese Literatur beziehen, gilt diese Einschränkung.

\subsection{Ansätze mit lokalen, kleinen Modellen}

Ein viertes Feld setzt lokal betreibbare Modelle geringer Größe ein, meist aus Gründen des Datenschutzes, der Latenz oder der Offline-Fähigkeit.
Ein Ansatz zur Handlungsplanung nutzt ein lokales Modell mit sieben Milliarden Parametern und zeigt, dass sich solche Aufgaben grundsätzlich auch mit kleinen Modellen bearbeiten lassen~\cite{wu2023}.
Ein weiterer löst das Auffinden von Objekten in 3D-Szenen mit einem kleinen visuellen Modell~\cite{dinh2026}.
Beide belegen die Machbarkeit unter Ressourcenbeschränkung, vergleichen aber nicht mehrere Repräsentationsformen für die Szenenmanipulation.
Welche Repräsentation eine natürlichsprachliche Manipulation unter diesen Bedingungen am zuverlässigsten trägt, ist damit offen.


% ------------------------------------------------------------
\section{Forschungslücke}
\label{sec:grundlagen-luecke}

Strukturelle und visuelle Repräsentation sind jeweils für sich etabliert.
Es gibt Arbeiten zur sprachgesteuerten Manipulation von Szenen, zum hybriden Auffinden von Objekten, zur Wirkung einzelner Modalitäten und zum Einsatz kleiner, lokal betreibbarer Modelle.
Keine bekannte Arbeit verbindet jedoch die folgenden vier Eigenschaften: Sie befasst sich mit der natürlichsprachlichen Manipulation einer Szene und nicht nur mit deren Verständnis, dem Auffinden von Objekten oder der Erzeugung von Szenen; sie vergleicht strukturelle, visuelle und hybride Repräsentation als gleichwertige Bedingungen; sie schlüsselt die Aufgaben nach der Art der sprachlichen Referenz auf; und sie setzt kleine, lokal betreibbare Modelle statt großer Cloud-Modelle ein.

Offen ist damit, wie eine Szene einem lokal betriebenen Agenten für die natürlichsprachliche Manipulation bereitgestellt werden sollte und welchen Einfluss die Repräsentationsform auf den Ausführungserfolg hat.
Insbesondere ist unklar, ob eine hybride Kombination die Stärken beider Repräsentationen vereint oder ob bei einem kleinen, lokalen Modell eine Quelle die andere verdrängt.
Hier setzen die beiden Forschungsfragen aus Abschnitt~\ref{sec:ein-ziele} an.
